\documentclass[10pt,a4paper]{amsart}
\usepackage[margin=19mm]{geometry}
\usepackage{amsmath,amssymb,mathtools}
\usepackage[scr=rsfs,cal=euler]{mathalfa}
\usepackage{xfrac}
\usepackage[unicode,hidelinks,bookmarks=false]{hyperref}
\newcommand{\ie}{i.e.\ }
\newcommand{\cf}{cf.\ }
\newcommand{\resp}{respectively\ }
\newcommand{\Subsection}{\subsection}
\providecommand{\nobreakdash}{\nobreak\hbox{-}}
\setlength{\emergencystretch}{2em}
\multlinegap0pt
\usepackage{amssymb}
\newtheorem{prop}{Proposition}
\newcommand{\R}{\mathbb{R}}
\title{Small Ball Probabilities for the Stochastic Heat Equation on Compact Manifolds}
\author{Jiaming Chen}
\date{Source: arXiv:2601.20794v1 (CC BY 4.0)}
\begin{document}
\maketitle

\noindent\emph{Source and licence.} Small Ball Probabilities for the Stochastic Heat Equation on Compact Manifolds. Version arXiv:2601.20794v1 (CC BY 4.0), distributed by arXiv under the Creative Commons Attribution 4.0 International licence, http://creativecommons.org/licenses/by/4.0/, as recorded in the arXiv licence metadata for this exact version and shown on the abstract page https://arxiv.org/abs/2601.20794v1. Full licence notice: \texttt{licenses/arxiv-2601.20794-CC-BY-4.0.txt}.\par\smallskip

\noindent\textit{Editorial scope.} Counted unit: the article's prop prop (source0.tex lines 302--318). The article's complete original proof of it is delivered with it (source0.tex lines 322--379, one \texttt{proof} environment of 58 lines, which the article places away from the statement). No mathematics is written, repaired or re-derived: the statement and the proof are the article's own lines, in the article's own notation, and no retained line is substituted or re-flowed.  Retained uncounted as the source-local context the proof needs: lines 73--76; lines 172--175; lines 284--287; lines 291--294.  Nothing was omitted from any retained span, so the package carries no omission record.  No retained line contains a citation, so the wrapper carries no bibliography and no bibliography entry.  No retained line contains a figure environment, an image-inclusion command or drawing code, so no image is delivered and the package declares no asset dependency.  Editorial formatting only, in addition: an AMS \texttt{amsart} preamble that restates the article's own declarations and macros used by the retained lines, namely the environments \texttt{prop} on separate counters, together with the standard packages those lines need.  The retained environments print in this document as Proposition 1 in order of appearance, whereas the article prints the same items with the numbering of its own full document.  The complete native source of the article as distributed in the arXiv e-print for this exact version is bundled unchanged as \texttt{sources/193549/source0.tex}.
\begin{equation}
\label{she}
\partial_t u(t,x)=\frac{1}{2}\Delta_M u(t,x)+\sigma(t,x,u)\dot{W}(t,x),\quad (t,x)\in \R_+\times M,
\end{equation}

\begin{equation}
\label{hypothesis1}
\vert \sigma(t,x,u)-\sigma(t,x,v)\vert\leq\mathcal{D}|u-v|.
\end{equation}

\begin{equation}
\label{Fn}
F_n=\left\lbrace\vert u(t,x)\vert\leq \sqrt{f(t_1)}\text{ for all } (t,x)\in R_{n,J}\right\rbrace,
\end{equation}

\begin{equation}
\label{En}
E_n=\left\lbrace\vert u(t_{n+1},x)\vert\leq \frac{\mathcal{C}_3}{3}t_1^{h/4},\,\text{and }\vert u(t,x)\vert\leq \mathcal{C}_3t_1^{h/4}\text{ for all } t\in[t_n,t_{n+1}),x\in M\right\rbrace.
\end{equation}

\begin{prop}
\label{prop}
Consider the solution to \eqref{she} with the initial condition $u_0\equiv 0$. There exist positive constants $\textbf{C}_4,\textbf{C}_5,\textbf{C}_6,\textbf{C}_7,\mathcal{D}_0$ and $\varepsilon_1$, independent of $\varepsilon$, such that for all $0<\varepsilon<\varepsilon_1$ and $0<\mathcal{D}<\mathcal{D}_0$ in \eqref{hypothesis1}, 
\begin{enumerate}
\item[(a)]
\begin{equation*}
P\left(F_n\bigg| \bigcap_{k=0}^{n-1}F_k\right)\leq\begin{cases}
    \textbf{C}_4\exp\left(-\frac{\textbf{C}_5}{\varepsilon^2}\right)&(d-1)/2<\alpha\leq d/2\\
    \textbf{C}_4\exp\left(-\frac{\textbf{C}_5}{t_1^{\min\left(\frac{d-2\alpha}{2d},\frac{2\alpha+2-d}{2}\right)}}\right)&\max(0,d/2-1)<\alpha\leq(d-1)/2,
\end{cases}
\end{equation*}
\item[(b)] 
\begin{equation*}
P\left(E_n\bigg| \bigcap_{k=-1}^{n-1}E_k\right)\geq \textbf{C}_6\exp\left(-\frac{\textbf{C}_7}{t_1^{d/2}}\right).
\end{equation*}
\end{enumerate}
\end{prop}

\begin{proof}
The event $F_{n}$ in $\eqref{Fn}$ controls the behavior of $u(t,x)$ at the discrete time $t_n$. Taking the intersection over all such times yields
\[F:=\bigcap_{n=0}^{\left\lfloor\frac{T}{t_1}\right\rfloor}F_{n}\supset \left\lbrace\vert u(t,x)\vert\leq \sqrt{f(t_1)}, t\in[0,T],x\in M\right\rbrace.\]
Moreover,
\[P(F)=P\left(\bigcap_{n=0}^{\left\lfloor\frac{T}{t_1}\right\rfloor}F_{n}\right)=P(F_{0})\prod_{n=1}^{\left\lfloor\frac{T}{t_1}\right\rfloor}P\left(F_{n}\bigg| \bigcap_{k=0}^{n-1}F_{k}\right).\]
Since $u_0(x)\equiv 0$, $F_0=\Omega$, and Proposition \ref{prop}(a) therefore implies that, for $\max(0,d/2-1)<\alpha\leq(d-1)/2$,
\begin{align*}
P\left(\left\lbrace\vert u(t,x)\vert\leq t_1^\frac{2\alpha+2-d}{4}, t\in[0,T],x\in M\right\rbrace\right)&\leq P(F)\leq \left[\textbf{C}_4\exp\left(-\frac{\textbf{C}_5}{t_1^{\min\left(\frac{d-2\alpha}{2d},\frac{2\alpha+2-d}{2}\right)}}\right)\right]^{\left\lfloor\frac{T}{t_1}\right\rfloor}\\
&\leq \textbf{C}_2\exp\left(-\frac{\textbf{C}_3T}{t_1^{\min\left(\frac{3d-2\alpha}{2d},\frac{2\alpha+4-d}{2}\right)}}\right).
\end{align*}
Replacing $t_1^\frac{2\alpha+2-d}{4}$ with $\varepsilon_2$ yields
\begin{equation*}
P\left(\left\lbrace\vert u(t,x)\vert\leq \varepsilon_2, t\in[0,T],x\in M\right\rbrace\right)\leq\textbf{C}_2\exp\left(-\frac{\textbf{C}_3T}{\varepsilon_2^{\min\left(\frac{6d-4\alpha}{d(2\alpha+2-d)},\frac{4\alpha+8-2d}{2\alpha+2-d}\right)}}\right).
\end{equation*}

For $(d-1)/2<\alpha< d/2$, Proposition \ref{prop}(a) gives
\[
P\left(\left\lbrace\vert u(t,x)\vert\leq t_1^\frac{2\alpha+2-d}{4}, t\in[0,T],x\in M\right\rbrace\right)\leq\textbf{C}_2\exp\left(-\frac{\textbf{C}_3T}{t_1^{\frac{3d-2\alpha}{2d}}}\right),
\]
and replacing $t_1^\frac{2\alpha+2-d}{4}$ with $\varepsilon_2$ yields
\begin{equation*}
P\left(\left\lbrace\vert u(t,x)\vert\leq \varepsilon_2, t\in[0,T],x\in M\right\rbrace\right)\leq\textbf{C}_2\exp\left(-\frac{\textbf{C}_3T}{\varepsilon_2^{\frac{6d-4\alpha}{d(2\alpha+2-d)}}}\right).
\end{equation*}

Note that the exponent $\frac{6d-4\alpha}{d(2\alpha+2-d)}$ is decreasing on $(d-1)/2<\alpha< d/2$, and
\begin{equation*}
    \inf_{\alpha}\frac{6d-4\alpha}{d(2\alpha+2-d)}=2,
\end{equation*}
as $\alpha\to \frac{d}{2}^-$. At the critical value $\alpha=d/2$, an additional logarithmic correction appears. Indeed,
\begin{equation*}
    P\left(\left\lbrace\vert u(t,x)\vert\leq \sqrt{t_1\ln(1/t_1)}, t\in[0,T],x\in M\right\rbrace\right)\leq\textbf{C}_2\exp\left(-\frac{\textbf{C}_3T}{t_1\varepsilon^2}\right).
\end{equation*}
Replacing $\sqrt{t_1\ln(1/t_1)}$ with $\varepsilon_2$, we obtain
\begin{equation*}
    t_1\varepsilon^2\leq\frac{C\varepsilon_2^2\varepsilon^{2+2d}}{\ln(1/\varepsilon)}\leq\frac{C\varepsilon_2^2\ln(1/\varepsilon_2)^{-C'}}{\ln\ln(1/\varepsilon_2)}\leq\frac{C\varepsilon_2^2}{[\ln\ln(1/\varepsilon_2)]^{2}},
\end{equation*}
which implies
\begin{equation*}
    P\left(\left\lbrace\vert u(t,x)\vert\leq \varepsilon_2, t\in[0,T],x\in M\right\rbrace\right)\leq\textbf{C}_2\exp\left(-\frac{\textbf{C}_3T[\ln\ln(1/\varepsilon_2)]^2}{\varepsilon_2^{2}}\right).
\end{equation*}

We now turn to the lower bound. The event $E_{n}$ in $\eqref{En}$ deals with the behavior of $u(t,x)$ over the entire interval $[t_n,t_{n+1}]$. Hence
\[E:=\bigcap_{n=-1}^{\left\lfloor\frac{T}{t_1}\right\rfloor}E_{n}\subset \left\lbrace\vert u(t,x)\vert\leq t_1^{h/4}, t\in[0,T],x\in M\right\rbrace,\]
and
\[P(E)=P\left(\bigcap_{n=-1}^{\left\lfloor\frac{T}{t_1}\right\rfloor}E_{n}\right)=P(E_{-1})\prod_{n=0}^{\left\lfloor\frac{T}{t_1}\right\rfloor}P\left(E_{n}\bigg| \bigcap_{k=-1}^{n-1}E_{k}\right).\]
With $u_0(x)\equiv 0$ and $E_{-1}=\Omega$, Proposition \ref{prop}(b) immediately yields
\begin{equation*}
    \begin{split}
        P\left(\left\lbrace\vert u(t,x)\vert\leq \mathcal{C}_3t_1^{h/4}, t\in[0,T],x\in M\right\rbrace\right)&\geq P(E)\geq \left[\textbf{C}_6\exp\left(-\frac{\textbf{C}_7}{t_1^{d/2}}\right)\right]^{\left\lfloor\frac{T}{t_1}\right\rfloor+1}\\
        &\geq \textbf{C}_0\exp\left(-\frac{\textbf{C}_1T}{t_1^{d/2+1}}\right).
    \end{split}
\end{equation*}
Finally, replacing $t_1^{h/4}$ with $\varepsilon_3$ gives
\begin{equation*}
P\left(\left\lbrace\vert u(t,x)\vert\leq \varepsilon_3, t\in[0,T],x\in M\right\rbrace\right)\geq\textbf{C}_0\exp\left(-\frac{\textbf{C}_1T}{\varepsilon_3^{\frac{2d+4}{h}}}\right).
\end{equation*}
The proof is complete, and the rest of this paper is devoted to the proof of Proposition \ref{prop}.
\end{proof}

\end{document}
